\documentclass[11pt, a4paper]{article}

\usepackage[a4paper, top=2.5cm, bottom=2.5cm, left=2cm, right=2cm]{geometry}
\usepackage{amsmath}
\usepackage{amssymb}
\usepackage[colorlinks=true, linkcolor=blue, urlcolor=blue, citecolor=blue]{hyperref}

\title{Theory of Everything - Volume 15 \\ \Large The Early Universe \& Cosmic Inflation}
\author{Omega Protocol}
\date{\today}

\begin{document}

\maketitle

\section*{Layman's Summary}
How did the universe begin? Volume 15 explores the very first moments of existence. We show that the universe did not start as a tiny, infinitely dense point of matter, but rather as a highly structured, purely informational state that rapidly "unzipped" to create the vast, expanding cosmos we see today.

\section{The Initial State (Axiom 18)}
We establish that the universe began at "Time Zero" in a state of maximal informational density but absolutely zero topological complexity. Imagine a tightly compressed ZIP file containing the blueprint for reality. It had no shape, no distance, and no geometry—just pure, highly ordered quantum information waiting to unfold.

\section{Cosmic Inflation (Theorem)}
In its first fraction of a second, the universe expanded at an unimaginably fast rate, a period known as Cosmic Inflation. We prove that this wasn't driven by a mysterious "inflaton field," but rather by a phase transition in the fundamental information network. The transition from a longitudinal informational state ($\Phi_N$) to a transverse state ($\Phi_\Delta$) violently expanded the geometry of space.

\section{Primordial Fluctuations (Theorem)}
If the initial state was perfectly smooth, how did we get galaxies and stars? We prove that tiny, inherent quantum jitters in the fundamental "modular flow" of the universe acted as seeds. As inflation stretched the universe, these microscopic informational vibrations were frozen into the macroscopic structure, eventually pulling matter together to form the cosmic web of galaxies.

\end{document}
